\section{Evaluation Protocol}\label{sec:exp:protocol}\label{sec:exp:results:protocol}

All scores in this chapter use the held-out test split of \cref{sec:exp:datasets}.
All models are scored at $L=2000$ with no overlap.
Neural models are scored on $1168$ windows from $244$ scenarios.
Feature \gls{dbscan} is not trained and is scored on $1412$ windows of the same length.
Clustering uses the $\varepsilon$ of \cref{tab:exp:hparams_attn}.

The reported scores are $\ARI$, $\AMI$, homogeneity, and completeness, defined in \cref{sec:method:metrics}, and the absolute cluster-count error $\lvert\hat{K}-K\rvert$ on the emitter branch and $\lvert\hat{M}-M\rvert$ on the mode branch.
Each table cell is the mean $\pm$ one standard deviation of that per-window score.

The next three sections answer the research questions of \cref{sec:intro:rq} in order.
\Cref{sec:exp:results} is the existence test for RQ1: a jointly trained Vanilla encoder against Feature \gls{dbscan}.
\Cref{sec:exp:ablation} isolates the joint objective (RQ2).
\Cref{sec:exp:attention} compares attention mechanisms under that joint objective (RQ3).
Qualitative plots of one input window and its embeddings follow in \cref{sec:exp:qualitative}.
